%% kapa: exam=1 round=33 year=2022 period=2 form=A subject=appraisal_law subject_ko=감정평가관계법규 seq=9 number=9
%% source: Q-Net 감정평가사 기출문제 (cst00309) artlSeq=5209997 file=raw/1008/5209997/2022년도 제33회 감정평가사 1차시험 2교시(감정평가 관계법규, 회계학).pdf
%% conversion: kordoc 4.24.51 -> tools/gen1.py (auto); review=proofread-sonnet; flags=none
\begin{question}{9}
국토의 계획 및 이용에 관한 법령상 A군수가 민간건설업자 B에 대해 개발행위허가를 할 때, 토석을 운반하는 차량 통행으로 통행로 주변 환경이 오염될 우려가 있어 환경오염 방지의 이행 보증 등에 관한 조치를 명하는 경우이다. 그에 관한 설명으로 옳은 것은?
\begin{choices}
  \item B가 예치하는 이행보증금은 총공사비의 30퍼센트 이상이 되도록 해야 한다.
  \item B가 준공검사를 받은 때에는 A군수는 즉시 이행보증금을 반환하여야 한다.
  \item A군수는 이행보증금을 행정대집행의 비용으로 사용할 수 없다.
  \item B가 산지에서 개발행위를 하는 경우 이행보증금의 예치금액의 기준이 되는 총공사비에는 「산지관리법」에 따른 복구비는 포함되지 않는다.
  \item B가 민간건설업자가 아닌 국가인 경우라도 민간건설업자의 경우와 동일한 이행보증이 필요하다.
\end{choices}
\end{question}
\answer{2}
